\documentclass[11pt,letterpaper]{article}

% ── Packages ──
\usepackage[margin=1in]{geometry}
\usepackage{amsmath,amssymb,amsthm,mathtools}
\usepackage{physics}
\usepackage{enumitem}
\usepackage{hyperref}
\usepackage{cleveref}
\usepackage{tcolorbox}
\usepackage{tikz}
\usetikzlibrary{arrows.meta,decorations.markings,calc,patterns}
\usepackage{pgfplots}
\pgfplotsset{compat=1.18}
\usepackage{fancyhdr}
\usepackage{titlesec}
\usepackage{booktabs}
\usepackage{cancel}
\usepackage{bm}

% ── Theorem environments ──
\theoremstyle{definition}
\newtheorem{definition}{Definition}[section]
\newtheorem{example}{Example}[section]
\theoremstyle{plain}
\newtheorem{theorem}{Theorem}[section]
\newtheorem{lemma}[theorem]{Lemma}
\newtheorem{proposition}[theorem]{Proposition}
\newtheorem{corollary}[theorem]{Corollary}
\theoremstyle{remark}
\newtheorem{remark}{Remark}[section]

% ── tcolorbox styles ──
\tcbuselibrary{skins,breakable}
\newtcolorbox{keyresult}[1][]{colback=blue!5!white,colframe=blue!60!black,
  fonttitle=\bfseries,title={#1},breakable}
\newtcolorbox{physicalnote}[1][]{colback=green!5!white,colframe=green!50!black,
  fonttitle=\bfseries,title={#1},breakable}
\newtcolorbox{warningbox}[1][]{colback=red!5!white,colframe=red!60!black,
  fonttitle=\bfseries,title={#1},breakable}

% ── Header/Footer ──
\pagestyle{fancy}
\fancyhf{}
\fancyhead[L]{\small Nonlinear Dynamics and Chaos}
\fancyhead[R]{\small\nouppercase{\leftmark}}
\fancyfoot[C]{\thepage}
\renewcommand{\headrulewidth}{0.4pt}

% ── Title formatting ──
\titleformat{\section}{\Large\bfseries}{\thesection}{1em}{}
\titleformat{\subsection}{\large\bfseries}{\thesubsection}{1em}{}

% ── Custom commands ──
\newcommand{\R}{\mathbb{R}}
\newcommand{\xdot}{\dot{x}}
\newcommand{\xstar}{x^*}
\newcommand{\thetadot}{\dot{\theta}}
\newcommand{\fp}[1]{f'(#1)}
\newcommand{\sgn}{\operatorname{sgn}}

\begin{document}

% ══════════════════════════════════════════════
% TITLE PAGE
% ══════════════════════════════════════════════
\begin{titlepage}
\centering
\vspace*{2cm}
{\Huge\bfseries Nonlinear Dynamics and Chaos\par}
\vspace{0.5cm}
{\Large Comprehensive Lecture Notes\par}
\vspace{1.5cm}
{\large Based on the course following\\[6pt]
\textit{Nonlinear Dynamics and Chaos} (3rd Edition)\\[4pt]
by Steven H.\ Strogatz\par}
\vspace{2cm}
{\large Spring 2026\par}
\vfill
{\small These notes are self-contained and assume familiarity with
multivariable calculus, linear algebra, and ordinary differential equations
at the level of a mathematical methods of physics course.\par}
\end{titlepage}

\tableofcontents
\newpage

% ══════════════════════════════════════════════
% PART I: ONE-DIMENSIONAL FLOWS
% ══════════════════════════════════════════════
\part{One-Dimensional Flows}

% ──────────────────────────────────────────────
\section{Flows on the Line}\label{sec:flows-line}
% ──────────────────────────────────────────────

\subsection{First-Order ODEs and Autonomous Systems}

The most general first-order ordinary differential equation (ODE) for a single
real variable $x(t)$ is
\begin{equation}\label{eq:general-1st}
  \xdot = f(x,t),
\end{equation}
where the overdot denotes $\dd{x}/\dd{t}$.

\begin{definition}[Autonomous ODE]
An ODE of the form
\begin{equation}\label{eq:autonomous}
  \xdot = f(x)
\end{equation}
is called \textbf{autonomous} because the right-hand side does not depend
explicitly on time $t$.  The function $f$ specifies the \emph{velocity}
of the state $x$ on the real line.
\end{definition}

The system~\eqref{eq:autonomous} is \textbf{linear} if $f(x)=ax+b$ for
constants $a,b$, and \textbf{nonlinear} otherwise.

Given an initial condition $x(0)=x_0$, the solution $x(t)$ traces a
trajectory on the real line.  Because $f$ depends only on $x$, the velocity
at a given position $x$ is the same regardless of when the particle passes
through $x$.  This is the hallmark of an autonomous system.

\subsection{A Geometric Way of Thinking: The Phase Line}

Rather than solving~\eqref{eq:autonomous} analytically (which is often
impossible for nonlinear $f$), we adopt a \emph{geometric} approach.

\medskip\noindent\textbf{Procedure.}
\begin{enumerate}[leftmargin=2em]
\item Plot $f(x)$ versus $x$.
\item Where $f(x)>0$, the flow moves to the \emph{right} (increasing $x$);
      where $f(x)<0$, the flow moves to the \emph{left}.
\item Points where $f(\xstar)=0$ are at rest---they are called
      \textbf{fixed points} (or equilibria, critical points, stationary
      points).
\end{enumerate}

\begin{definition}[Fixed Point]
A point $\xstar$ is a \textbf{fixed point} of $\xdot = f(x)$ if
$f(\xstar)=0$.
\end{definition}

\begin{definition}[Stability of a Fixed Point---Informal]
A fixed point $\xstar$ is
\begin{itemize}[leftmargin=2em]
\item \textbf{stable} (an ``attractor'') if nearby trajectories approach
      $\xstar$ as $t\to\infty$;
\item \textbf{unstable} (a ``repeller'') if nearby trajectories move away
      from $\xstar$.
\end{itemize}
\end{definition}

\noindent\textbf{Graphical rule.}
Looking at the graph of $f(x)$: a fixed point where $f$ crosses from
positive to negative (i.e., $\fp{\xstar}<0$) is stable; a crossing from
negative to positive ($\fp{\xstar}>0$) is unstable.

\begin{physicalnote}[Direction Field]
For autonomous systems, the direction field has identical slope $f(x)$ at
every point along a horizontal line $x = \text{const}$.  Solution
trajectories in the $(t,x)$ plane are parallel translations of each other.
The initial condition $x_0$ simply selects which trajectory the system
follows.
\end{physicalnote}

\subsection{Linear Stability Analysis}\label{sec:lin-stab-1d}

To make the stability criterion precise, we linearize near the fixed point.

Let $\xstar$ be a fixed point and set $u = x - \xstar$ (a small
perturbation).  Then
\[
  \dot{u} = \xdot = f(x) = f(\xstar) + \fp{\xstar}\,u + O(u^2).
\]
Since $f(\xstar)=0$, dropping higher-order terms gives the \emph{linearized
equation}:
\begin{equation}\label{eq:lin-1d}
  \dot{u} \approx \fp{\xstar}\, u.
\end{equation}
This is a linear ODE with solution
\begin{equation}\label{eq:lin-1d-soln}
  u(t) = u_0 \, e^{\fp{\xstar}\, t}.
\end{equation}

\begin{keyresult}[Linear Stability Criterion (1D)]
Let $\xstar$ be a fixed point of $\xdot = f(x)$.
\begin{itemize}[leftmargin=2em]
\item If $\fp{\xstar} < 0$, then $\xstar$ is \textbf{linearly stable}.
\item If $\fp{\xstar} > 0$, then $\xstar$ is \textbf{unstable}.
\item If $\fp{\xstar} = 0$, the test is \textbf{inconclusive}---higher-order
      terms must be examined.
\end{itemize}
\end{keyresult}

\begin{example}[Population Growth]
Consider $\xdot = rx(1-x/K)$ (the logistic equation) with $r,K>0$.
The fixed points are $\xstar=0$ and $\xstar=K$.
Since $f'(x) = r - 2rx/K$, we have $f'(0)=r>0$ (unstable) and
$f'(K)=-r<0$ (stable).  The carrying capacity $K$ is the stable
equilibrium.
\end{example}

\subsection{Existence and Uniqueness}\label{sec:existence-uniqueness}

\begin{theorem}[Existence and Uniqueness]\label{thm:exist-unique}
Consider the initial value problem $\xdot = f(x)$, $x(0)=x_0$.
If $f(x)$ is continuous and $f'(x)$ is also continuous (i.e.,
$f\in C^1$) on an open interval containing $x_0$, then a solution $x(t)$
exists and is unique on some time interval $(-\tau,\tau)$ about $t=0$.
\end{theorem}

\begin{remark}
A key consequence is that trajectories on the real line \emph{cannot cross}
fixed points.  This means the qualitative behavior is completely determined
by the signs of $f(x)$ between consecutive fixed points.
\end{remark}

\subsection{Impossibility of Oscillations}

\begin{theorem}
For a first-order autonomous ODE $\xdot=f(x)$ on the real line with $f\in C^1$, oscillations (periodic solutions) are impossible.
\end{theorem}
\begin{proof}
The velocity $f(x)$ at any point $x$ is single-valued and time-independent.  For the system to oscillate, $x$ would have to return to a previously visited value, requiring $\xdot$ to change sign while passing through the same point---but $f(x)$ is fixed at each $x$.  More precisely, by the uniqueness theorem, trajectories cannot cross, and on $\R^1$ there is no way to return to a starting point without crossing a fixed point, at which the solution would be trapped forever.
\end{proof}

\subsection{Potentials}

For one-dimensional flows, we can often think of the dynamics in terms of a ``particle rolling downhill on a potential landscape.''

\begin{definition}[Potential Function]
Given $\xdot = f(x)$, define the \textbf{potential} $V(x)$ by
\begin{equation}
  f(x) = -\dv{V}{x}, \qquad\text{i.e.}\qquad V(x) = -\int f(x)\,dx.
\end{equation}
\end{definition}

Then $\xdot = -V'(x)$: the system always moves in the direction of
decreasing $V$.  Stable fixed points correspond to local \emph{minima} of
$V$; unstable fixed points correspond to local \emph{maxima}.

The potential decreases monotonically along trajectories:
\[
  \dv{V}{t} = V'(x)\,\xdot = V'(x)\bigl[-V'(x)\bigr]
  = -\bigl[V'(x)\bigr]^2 \le 0.
\]

\subsection{Solving Equations on the Computer: Euler's Method}

In practice, many nonlinear ODEs cannot be solved in closed form.
The simplest numerical scheme for $\xdot = f(x)$ is \textbf{Euler's method}:
\begin{equation}
  x_{n+1} = x_n + f(x_n)\,\Delta t,
\end{equation}
which approximates the solution by taking small steps of size $\Delta t$
along the tangent line.  More sophisticated methods (Runge--Kutta, adaptive
step-size) are used in practice, but Euler's method captures the essential
idea.

\begin{example}[Separation of Variables and Finite-Time Blowup]
Consider $\xdot = 1 + x^2$ with $x(0)=0$.  By separation of variables:
\[
  \int\frac{dx}{1+x^2} = \int dt
  \quad\Longrightarrow\quad \arctan(x) = t + C.
\]
With $x(0)=0$: $C=0$, so $x(t) = \tan(t)$.  The solution blows up in
finite time $t = \pi/2$.

This illustrates an important caveat: the existence and uniqueness theorem
guarantees solutions only \emph{locally} in time.  If $f$ grows fast enough,
the solution can escape to infinity in finite time.  However, for flows on
the \emph{line} where $f$ is smooth and bounded, this cannot happen.
\end{example}

\begin{example}[Exact Solution: $\xdot = x^2 - 1$]
Fixed points: $\xstar = \pm 1$.  Since $f'(x)=2x$, $f'(-1)=-2<0$ (stable)
and $f'(1)=2>0$ (unstable).  By partial fractions:
\[
  \int\frac{dx}{x^2-1} = \int\frac{dx}{(x-1)(x+1)}
  = \frac{1}{2}\int\left(\frac{1}{x-1} - \frac{1}{x+1}\right)dx = t + C.
\]
Solving gives $x(t) = -\tanh(t-t_0)$ for trajectories between the fixed
points, confirming the graphical analysis.
\end{example}


% ──────────────────────────────────────────────
\section{Bifurcations}\label{sec:bifurcations}
% ──────────────────────────────────────────────

\subsection{Introduction}

A \textbf{bifurcation} occurs when a small, smooth change in a parameter
causes a qualitative (topological) change in the dynamics---for instance,
fixed points may be created, destroyed, or change stability.

The parameter is traditionally denoted $r$.  The key question is:
\emph{How does the phase portrait of $\xdot = f(x;r)$ change as $r$ varies?}

A \textbf{bifurcation diagram} plots a characteristic feature of the
solution (usually the fixed point values $\xstar$) on the vertical axis
versus the parameter $r$ on the horizontal axis.  Solid curves denote
stable fixed points; dashed curves denote unstable ones.

\subsection{Saddle-Node Bifurcation}\label{sec:saddle-node}

The \textbf{saddle-node bifurcation} is the basic mechanism by which fixed
points are created and destroyed.

\begin{keyresult}[Normal Form: Saddle-Node Bifurcation]
\begin{equation}\label{eq:sn-normal}
  \xdot = r + x^2.
\end{equation}
\end{keyresult}

\noindent\textbf{Analysis.}
\begin{itemize}[leftmargin=2em]
\item For $r<0$: two fixed points at $\xstar = \pm\sqrt{-r}$.
      Since $f'(x)=2x$, the negative root is stable ($f'<0$) and the
      positive root is unstable ($f'>0$).
\item At $r=0$: the two fixed points coalesce into one \emph{half-stable}
      fixed point at $\xstar=0$.
\item For $r>0$: no fixed points exist---$f(x)>0$ everywhere so $x$
      increases without bound.
\end{itemize}

The bifurcation occurs at $r_c = 0$.  Graphically, the two branches of the
bifurcation diagram meet in a cusp.

\begin{physicalnote}[Bottleneck / Ghost Effect]
Just above the bifurcation ($r = 0^+$), even though no fixed point exists,
the system slows dramatically near $x=0$ because $f(x)\approx r$ is small.
The time to pass through this ``bottleneck'' scales as
$T_{\text{bottleneck}} \sim \pi/\sqrt{r}$ for $r\to 0^+$.
This is sometimes called the \emph{ghost} of the vanished fixed point.
\end{physicalnote}

\subsection{Transcritical Bifurcation}\label{sec:transcritical}

In many physical situations, a fixed point exists for all parameter values
but exchanges stability with another fixed point.

\begin{keyresult}[Normal Form: Transcritical Bifurcation]
\begin{equation}\label{eq:tc-normal}
  \xdot = rx - x^2 = x(r-x).
\end{equation}
\end{keyresult}

\noindent\textbf{Analysis.}  Fixed points: $\xstar = 0$ and $\xstar = r$.
Since $f'(x) = r - 2x$:
\begin{itemize}[leftmargin=2em]
\item $f'(0) = r$: origin is stable for $r<0$, unstable for $r>0$.
\item $f'(r) = -r$: the fixed point $\xstar=r$ is unstable for $r<0$,
      stable for $r>0$.
\end{itemize}

At $r=0$ the two fixed points collide and \textbf{exchange stability}.
Neither fixed point is destroyed; they simply swap roles.

\begin{physicalnote}[Laser Threshold]
A laser below threshold ($r<0$) has zero output (stable $\xstar=0$).
Above threshold ($r>0$), $\xstar=0$ becomes unstable and a new stable
state $\xstar=r$ (proportional to the laser output intensity) takes over.
This is a transcritical bifurcation.
\end{physicalnote}

\subsection{Pitchfork Bifurcation}\label{sec:pitchfork}

Pitchfork bifurcations arise in systems that possess a \textbf{symmetry}
$x\to -x$.  There are two types.

\subsubsection{Supercritical Pitchfork}

\begin{keyresult}[Normal Form: Supercritical Pitchfork]
\begin{equation}\label{eq:pf-super}
  \xdot = rx - x^3.
\end{equation}
\end{keyresult}

\noindent\textbf{Analysis.}  Fixed points: $\xstar = 0$ for all $r$, and
$\xstar = \pm\sqrt{r}$ for $r>0$.

Since $f'(x) = r - 3x^2$:
\begin{itemize}[leftmargin=2em]
\item $f'(0) = r$: origin is stable for $r<0$, unstable for $r>0$.
\item $f'(\pm\sqrt{r}) = r - 3r = -2r < 0$ for $r>0$: the two new
      branches are stable.
\end{itemize}

At $r=0$, the origin loses stability and two new stable branches are born
symmetrically.  The bifurcation diagram looks like a ``pitchfork.''

\subsubsection{Subcritical Pitchfork}

\begin{keyresult}[Normal Form: Subcritical Pitchfork]
\begin{equation}\label{eq:pf-sub}
  \xdot = rx + x^3.
\end{equation}
\end{keyresult}

Here the cubic term is \emph{destabilizing}.  For $r<0$, there are three
fixed points: the origin (stable) and $\xstar = \pm\sqrt{-r}$ (both unstable).  At $r=0$ the unstable branches merge with the origin, making it
unstable for $r>0$.

\begin{warningbox}[Subcritical Pitchfork and Hysteresis]
In a physical system, the subcritical pitchfork is often ``stabilized'' by
a higher-order term, e.g., $\xdot = rx + x^3 - x^5$.  This leads to
\textbf{hysteresis}: the system can jump discontinuously between branches as
$r$ is varied, and the value of $x$ observed depends on the
\emph{direction} in which $r$ is changed.
\end{warningbox}

\subsection{Imperfect Bifurcations and Catastrophes}

When the symmetry of a pitchfork bifurcation is slightly broken (an
``imperfection'' $h$), the bifurcation diagram deforms.  The normal form is
\begin{equation}
  \xdot = h + rx - x^3.
\end{equation}
For $h\neq 0$, the perfect pitchfork breaks into an isolated branch and an
$S$-shaped curve.  The set of bifurcation points in the $(r,h)$ parameter
plane forms a cusp---this is the \textbf{cusp catastrophe}, one of the
elementary catastrophes of singularity theory.

\subsection{Summary of 1D Bifurcations}

\begin{center}
\begin{tabular}{@{}lcc@{}}
\toprule
\textbf{Bifurcation} & \textbf{Normal Form} & \textbf{Mechanism} \\
\midrule
Saddle-node     & $\xdot = r + x^2$       & Two FPs created/destroyed \\
Transcritical   & $\xdot = rx - x^2$      & FPs exchange stability \\
Supercrit.\ pitchfork & $\xdot = rx - x^3$ & One FP $\to$ three; stable branches \\
Subcrit.\ pitchfork   & $\xdot = rx + x^3$ & Three FPs $\to$ one; unstable branches \\
\bottomrule
\end{tabular}
\end{center}

\subsection{Application: Overdamped Bead on a Rotating Hoop}\label{sec:bead-hoop}

A bead slides without friction on a circular wire hoop of radius $R$ that
rotates about its vertical diameter with angular velocity $\omega$.  In the
overdamped limit (viscous damping dominates inertia), the angle $\phi$ of
the bead from the bottom of the hoop obeys:
\begin{equation}
  \dot{\phi} = f(\phi) = \sin\phi\bigl(\gamma\cos\phi - 1\bigr),
\end{equation}
where $\gamma = R\omega^2/g$ is the dimensionless rotation rate.

\textbf{Analysis:}
\begin{itemize}[leftmargin=2em]
\item Fixed points: $\phi^* = 0,\,\pi$ always exist.  For $\gamma > 1$,
      two additional fixed points appear at $\phi^* = \pm\arccos(1/\gamma)$.
\item At $\gamma = 1$: a \textbf{supercritical pitchfork bifurcation}
      occurs.  The bottom ($\phi=0$) changes from stable to unstable,
      and the bead spontaneously climbs to one of the new off-axis
      equilibria.
\end{itemize}

This is a beautiful physical realization of symmetry-breaking: for slow
rotation the bead sits at the bottom; for fast rotation it ``chooses'' a
side, breaking the left-right symmetry.

\subsection{Application: Insect Outbreak (Spruce Budworm)}

The budworm population $N$ satisfies a model of the form
\begin{equation}
  \dot{N} = rN\!\left(1 - \frac{N}{K}\right) - p(N),
\end{equation}
where the first term is logistic growth and $p(N) = BN^2/(A^2+N^2)$ is
predation (a sigmoidal functional response).  Nondimensionalizing
with $x = N/A$, $\tau = rAt$, etc., leads to
\begin{equation}
  \dot{x} = rx\!\left(1 - \frac{x}{k}\right) - \frac{x^2}{1+x^2}.
\end{equation}

For appropriate parameter ranges, this system exhibits \textbf{hysteresis}
via saddle-node bifurcations.  As parameters change slowly (e.g., the
forest matures, increasing $k$), the budworm population can undergo a
sudden catastrophic outbreak (a jump to a high-$x$ branch), followed by
an equally sudden crash when the branch disappears.  This involves two
saddle-node bifurcations forming a hysteresis loop---an example of a
\textbf{cusp catastrophe} in an ecological setting.


% ──────────────────────────────────────────────
\section{Flows on the Circle}\label{sec:flows-circle}
% ──────────────────────────────────────────────

\subsection{Motivation and Definitions}

On the real line, oscillations are impossible for a first-order autonomous
system.  However, if the phase space is a \textbf{circle} $S^1$
(i.e., the variable $\theta$ is $2\pi$-periodic), oscillations \emph{are}
possible: the system can flow around the circle indefinitely.

\begin{definition}
A \textbf{flow on the circle} is an ODE of the form
\begin{equation}
  \thetadot = f(\theta), \qquad \theta \in [0, 2\pi),
\end{equation}
where $f(\theta+2\pi) = f(\theta)$.
\end{definition}

Fixed points are again defined by $f(\theta^*)=0$, and stability analysis is
identical to the real-line case.

\subsection{Uniform and Nonuniform Oscillators}

\textbf{Uniform oscillator:} $\thetadot = \omega$ (constant angular
velocity).  Period $T = 2\pi/\omega$.

\textbf{Nonuniform oscillator:} $\thetadot = \omega - a\sin\theta$
with $a,\omega > 0$.
\begin{itemize}[leftmargin=2em]
\item If $\omega > a$: $\thetadot > 0$ always, so $\theta$ oscillates but
      nonuniformly.  The period is
      \begin{equation}\label{eq:nonuniform-period}
        T = \frac{2\pi}{\sqrt{\omega^2 - a^2}}.
      \end{equation}
\item At $\omega = a$: a saddle-node bifurcation creates a half-stable
      fixed point.  The period diverges: $T\to\infty$.
\item For $\omega < a$: two fixed points appear (one stable, one unstable)
      and oscillation ceases.
\end{itemize}

This is the prototypical example of a \textbf{saddle-node bifurcation on a
circle}, also called a SNIPER (Saddle-Node Infinite-PERiod) bifurcation.
The divergence of the period scales as $T \sim 1/\sqrt{\omega-a}$ near the
bifurcation.

\subsection{The Overdamped Pendulum}

A pendulum in a viscous fluid with a constant applied torque $\Gamma$ is
modeled by
\begin{equation}
  \thetadot = \Gamma - \sin\theta.
\end{equation}
This is exactly the nonuniform oscillator with $\omega = \Gamma$ and $a=1$.
For $|\Gamma|<1$, the pendulum settles to a tilted equilibrium; for
$|\Gamma|>1$, it whirls continuously.

\subsection{Superconducting Josephson Junctions}\label{sec:josephson-1d}

A Josephson junction consists of two superconductors separated by a thin
insulating barrier.  The \emph{phase difference} $\phi$ across the junction
obeys (in the overdamped, zero-capacitance limit):
\begin{equation}\label{eq:JJ-1D}
  \frac{\hbar}{2eR}\,\dot{\phi} = I - I_c \sin\phi,
\end{equation}
where $I$ is the applied (DC) current, $I_c$ is the critical current, and
$R$ is the shunt resistance.  Nondimensionalizing:
\begin{equation}
  \dot{\phi} = i - \sin\phi, \qquad i \equiv I/I_c.
\end{equation}
This is a flow on the circle identical to the overdamped pendulum.
\begin{itemize}[leftmargin=2em]
\item For $|i|<1$: two fixed points exist; $\dot{\phi}=0$ at steady state
      and the voltage $V = (\hbar/2e)\,\dot{\phi} = 0$.  The junction
      carries a \textbf{supercurrent} with zero resistance.
\item For $|i|>1$: no fixed points; $\phi$ winds continuously and a
      nonzero DC voltage appears across the junction.
\end{itemize}

% ══════════════════════════════════════════════
\part{Two-Dimensional Flows}
% ══════════════════════════════════════════════

% ──────────────────────────────────────────────
\section{Linear Systems in Two Dimensions}\label{sec:linear-2d}
% ──────────────────────────────────────────────

\subsection{Setup}

We now consider systems of two coupled first-order ODEs:
\begin{equation}\label{eq:2d-linear}
  \dot{\vec{x}} = A\vec{x}, \qquad
  A = \begin{pmatrix} a & b \\ c & d \end{pmatrix}, \qquad
  \vec{x} = \begin{pmatrix} x \\ y \end{pmatrix}.
\end{equation}

\subsection{Solution by Eigenvalues and Eigenvectors}

We seek solutions of the form $\vec{x}(t) = e^{\lambda t}\vec{v}$.
Substituting into~\eqref{eq:2d-linear}:
\[
  \lambda e^{\lambda t}\vec{v} = A\, e^{\lambda t}\vec{v}
  \quad\Longrightarrow\quad A\vec{v} = \lambda\vec{v}.
\]
Thus $\lambda$ must be an \textbf{eigenvalue} of $A$ and $\vec{v}$ the
corresponding \textbf{eigenvector}.

The eigenvalues satisfy the characteristic equation:
\begin{equation}\label{eq:char-eq}
  \lambda^2 - \tau\lambda + \Delta = 0,
\end{equation}
where $\tau \equiv \tr A = a+d$ and $\Delta \equiv \det A = ad-bc$.
The solutions are:
\begin{equation}\label{eq:eigenvalues}
  \lambda_{1,2} = \frac{\tau \pm \sqrt{\tau^2 - 4\Delta}}{2}.
\end{equation}

The general solution is $\vec{x}(t) = c_1 e^{\lambda_1 t}\vec{v}_1
+ c_2 e^{\lambda_2 t}\vec{v}_2$ (when the eigenvalues are distinct).

\subsection{Classification of Fixed Points}\label{sec:classification}

The qualitative behavior of trajectories near the origin depends on the
eigenvalues, which are determined by $\tau$ and $\Delta$.

\begin{keyresult}[Classification in the $(\tau,\Delta)$ Plane]
\begin{center}
\renewcommand{\arraystretch}{1.3}
\begin{tabular}{@{}lll@{}}
\toprule
\textbf{Region} & \textbf{Eigenvalues} & \textbf{Fixed Point Type} \\
\midrule
$\Delta < 0$ & Real, opposite signs & \textbf{Saddle point} \\
$\Delta > 0$, $\tau^2 > 4\Delta$, $\tau < 0$ & Real, both negative
  & \textbf{Stable node} \\
$\Delta > 0$, $\tau^2 > 4\Delta$, $\tau > 0$ & Real, both positive
  & \textbf{Unstable node} \\
$\Delta > 0$, $\tau^2 < 4\Delta$, $\tau < 0$ & Complex, $\text{Re}<0$
  & \textbf{Stable spiral} \\
$\Delta > 0$, $\tau^2 < 4\Delta$, $\tau > 0$ & Complex, $\text{Re}>0$
  & \textbf{Unstable spiral} \\
$\Delta > 0$, $\tau = 0$ & Pure imaginary
  & \textbf{Center} (linear) \\
$\Delta > 0$, $\tau^2 = 4\Delta$ & Real, repeated
  & \textbf{Star node} or \textbf{degenerate node} \\
$\Delta = 0$ & One eigenvalue zero
  & \textbf{Line of fixed points} \\
\bottomrule
\end{tabular}
\end{center}
\end{keyresult}

\noindent\textbf{Key relationships:}
\begin{align}
  \lambda_1 + \lambda_2 &= \tau = \tr A, \\
  \lambda_1 \cdot \lambda_2 &= \Delta = \det A.
\end{align}
The stability boundary is $\tau = 0$ (eigenvalues cross the imaginary axis).
The boundary between nodes and spirals is $\tau^2 = 4\Delta$.
Saddle points live below the $\Delta = 0$ axis.

\begin{physicalnote}[Physical Interpretation]
\begin{itemize}[leftmargin=2em]
\item \textbf{Stable node:} trajectories approach the origin along
      straight lines (the eigenvector directions); overdamped approach.
\item \textbf{Stable spiral:} trajectories spiral inward; underdamped
      oscillatory decay.
\item \textbf{Center:} closed elliptical orbits; no damping, no growth.
\item \textbf{Saddle:} trajectories approach along one eigenvector
      direction and are repelled along the other.
\end{itemize}
\end{physicalnote}


% ──────────────────────────────────────────────
\section{Phase Plane Analysis}\label{sec:phase-plane}
% ──────────────────────────────────────────────

\subsection{Nonlinear Systems and Phase Portraits}

We now consider the general two-dimensional nonlinear system:
\begin{equation}\label{eq:2d-nonlinear}
  \dot{x} = f(x,y), \qquad \dot{y} = g(x,y).
\end{equation}
The \textbf{phase portrait} is the collection of all trajectories in the
$(x,y)$ plane.

\begin{definition}[Nullclines]
The \textbf{$x$-nullcline} is the set of points where $\dot{x}=0$, i.e.,
$f(x,y)=0$.  The \textbf{$y$-nullcline} is the set where $\dot{y}=0$, i.e.,
$g(x,y)=0$.  Fixed points lie at the \emph{intersections} of the two
nullclines.
\end{definition}

Nullclines are extremely useful for sketching phase portraits:
on the $x$-nullcline, the flow is purely vertical ($\dot{x}=0$); on the
$y$-nullcline, the flow is purely horizontal ($\dot{y}=0$).

\subsection{Existence, Uniqueness, and Topological Consequences}

\begin{theorem}[Existence and Uniqueness in 2D]
If $f$ and $g$ are continuously differentiable, then given any initial
condition $(x_0,y_0)$, a unique solution exists (at least locally in time).
\end{theorem}

\begin{corollary}
Different trajectories can never cross in the phase plane.
\end{corollary}

This has profound consequences: for example, a closed orbit divides the plane
into an ``inside'' and an ``outside,'' and no trajectory can cross it.

\subsection{Linearization Near Fixed Points}\label{sec:linearization}

Let $(\xstar, y^*)$ be a fixed point of~\eqref{eq:2d-nonlinear}.  Set
$u = x - \xstar$, $v = y - y^*$.  Taylor expanding:
\begin{equation}
  \begin{pmatrix} \dot{u} \\ \dot{v} \end{pmatrix}
  = J\bigg|_{(\xstar,y^*)} \begin{pmatrix} u \\ v \end{pmatrix}
  + \text{h.o.t.},
\end{equation}
where $J$ is the \textbf{Jacobian matrix}:
\begin{equation}\label{eq:jacobian}
  J = \begin{pmatrix}
    \pdv{f}{x} & \pdv{f}{y} \\[4pt]
    \pdv{g}{x} & \pdv{g}{y}
  \end{pmatrix}\bigg|_{(\xstar,y^*)}.
\end{equation}

\begin{theorem}[Hartman--Grobman Theorem---Informal]\label{thm:hartman-grobman}
If the fixed point is \textbf{hyperbolic} (i.e., $\text{Re}(\lambda_i)\neq 0$
for all eigenvalues of $J$), then the phase portrait of the nonlinear system
near the fixed point is \emph{qualitatively the same} as the phase portrait
of the linearized system $\dot{\vec{u}} = J\vec{u}$.
\end{theorem}

\begin{warningbox}[Borderline Cases]
\begin{itemize}[leftmargin=2em]
\item \textbf{Centers} ($\tau=0$, $\Delta>0$) are \emph{not} hyperbolic.
      A linear center may become a spiral (stable or unstable) in the
      nonlinear system.  Additional analysis (e.g., conserved quantities,
      reversibility) is needed.
\item \textbf{Non-isolated fixed points} ($\Delta=0$) also require
      special treatment.
\end{itemize}
\end{warningbox}

\subsection{Application: Competing Species (Rabbits vs.\ Sheep)}

Consider the Lotka--Volterra competition model:
\begin{equation}
  \dot{x} = x(3-x-2y), \qquad \dot{y} = y(2-x-y),
\end{equation}
where $x$ and $y$ represent two competing populations.

\textbf{Nullclines:}
$\dot{x}=0$: $x=0$ or $x+2y=3$.  $\dot{y}=0$: $y=0$ or $x+y=2$.

\textbf{Fixed points:} $(0,0)$, $(3,0)$, $(0,2)$, and the intersection of
$x+2y=3$ and $x+y=2$, giving $(1,1)$.

Linearizing at each fixed point using the Jacobian:
\begin{itemize}[leftmargin=2em]
\item $(0,0)$: $J = \bigl(\begin{smallmatrix} 3&0\\0&2\end{smallmatrix}\bigr)$,
      eigenvalues $3,2 > 0$: unstable node.
\item $(3,0)$: $J = \bigl(\begin{smallmatrix} -3&-6\\0&-1\end{smallmatrix}\bigr)$,
      eigenvalues $-3,-1$: stable node (sheep extinct).
\item $(0,2)$: $J = \bigl(\begin{smallmatrix} -1&0\\-2&-2\end{smallmatrix}\bigr)$,
      eigenvalues $-1,-2$: stable node (rabbits extinct).
\item $(1,1)$: $\tau = -3$, $\Delta = 1$, $\tau^2-4\Delta = 5 > 0$:
      a saddle point (coexistence equilibrium is unstable).
\end{itemize}

The separatrix of the saddle at $(1,1)$ divides the first quadrant into two
basins of attraction: one leading to $(3,0)$ (species $x$ wins) and the
other to $(0,2)$ (species $y$ wins).  The outcome depends on initial conditions---\textbf{competitive exclusion}.

\subsection{Conservative Systems}\label{sec:conservative}

\begin{definition}[Conservative System]
A system $\dot{x}=f(x,y)$, $\dot{y}=g(x,y)$ is \textbf{conservative} if
there exists a function $E(x,y)$ (the \emph{energy} or \emph{conserved
quantity}) such that $E$ is constant along trajectories:
$\dv{E}{t} = 0$.
\end{definition}

\begin{theorem}
Conservative systems have \textbf{no attractors}.  Trajectories lie on
level sets of $E$, and there can be no stable nodes, stable spirals, or
stable limit cycles.
\end{theorem}
\begin{proof}
Since $E$ is constant along trajectories, no trajectory can approach a
single point (which would require $E$ to take a single value over an open
neighborhood), nor can a trajectory spiral inward (which would require $E$ to
decrease monotonically).
\end{proof}

\begin{example}[Harmonic Oscillator]
$m\ddot{x} = -kx$ can be written as
$\dot{x} = v$, $\dot{v} = -(k/m)\,x$.  The energy
$E = \tfrac{1}{2}kx^2 + \tfrac{1}{2}mv^2$
is conserved:
\[
  \dv{E}{t} = kx\dot{x} + mv\dot{v} = kxv + mv\bigl(-\tfrac{k}{m}x\bigr) = 0.
\]
The Jacobian at the origin has $\tau=0$, $\Delta=k/m>0$: a linear center.
Since the system is conservative, this center persists as a true nonlinear
center with closed elliptical orbits.
\end{example}

\subsection{Reversible Systems}\label{sec:reversible}

\begin{definition}[Time-Reversibility]
A system is \textbf{reversible} if it is invariant under the transformation
$t\to -t$ and some transformation of the state variables (typically
$y\to -y$, or equivalently, the velocity reverses).
\end{definition}

\begin{theorem}
For a reversible system, if a trajectory lies on the line of symmetry (the
axis that is invariant under the reversal), then any linear center is a
\emph{true} nonlinear center.
\end{theorem}

\begin{example}[Undamped Pendulum]
The equation $\ddot{\theta} = -\sin\theta$ written as a system:
$\dot{\theta} = v$, $\dot{v} = -\sin\theta$.
Under $t\to -t$, $v\to -v$: the equations are unchanged.  Therefore the
linear center at $(\theta,v)=(0,0)$ is a true nonlinear center.

The energy is $E = \tfrac{1}{2}v^2 - \cos\theta$.  Phase portrait:
\begin{itemize}[leftmargin=2em]
\item $E < 1$: librations (closed orbits around the center).
\item $E = 1$: the \textbf{separatrix} (homoclinic orbits connecting the
      saddle at $(\pi,0)$ to itself).
\item $E > 1$: rotations (unbounded trajectories wrapping around).
\end{itemize}
\end{example}

\subsection{Index Theory}\label{sec:index-theory}

Index theory provides topological constraints on the arrangement of fixed
points and closed orbits.

\begin{definition}[Index of a Closed Curve]
Let $C$ be a simple closed curve in the phase plane that does not pass
through any fixed points.  At each point on $C$, the vector field
$(\dot{x},\dot{y})$ makes an angle $\phi$ with the positive $x$-axis.
As we traverse $C$ counterclockwise, $\phi$ changes continuously.  The
\textbf{index} of $C$ is
\begin{equation}
  I_C = \frac{1}{2\pi}\oint_C d\phi = \frac{[\phi]_C}{2\pi},
\end{equation}
i.e., the number of complete counterclockwise revolutions of the vector
field as we traverse $C$ once counterclockwise.
\end{definition}

\begin{theorem}[Properties of the Index]
\begin{enumerate}[leftmargin=2em]
\item If $C$ encloses no fixed points, then $I_C = 0$.
\item If $C$ is a trajectory (e.g., a closed orbit), then $I_C = +1$.
\item The index is unchanged if $C$ is continuously deformed, as long as
      it does not cross a fixed point.
\item The index of a curve enclosing multiple isolated fixed points equals
      the sum of the indices of the individual fixed points.
\item Reversing time ($t\to -t$) does not change the index.
\end{enumerate}
\end{theorem}

\begin{keyresult}[Indices of Common Fixed Points]
\begin{center}
\begin{tabular}{@{}lc@{}}
\toprule
\textbf{Fixed Point Type} & \textbf{Index} \\
\midrule
Stable/unstable node    & $+1$ \\
Stable/unstable spiral  & $+1$ \\
Center                  & $+1$ \\
Saddle point            & $-1$ \\
\bottomrule
\end{tabular}
\end{center}
\end{keyresult}

\begin{corollary}
Any closed orbit must enclose fixed points whose indices sum to $+1$.
In particular, a closed orbit must enclose at least one fixed point that
is not a saddle.
\end{corollary}


% ──────────────────────────────────────────────
\section{Limit Cycles}\label{sec:limit-cycles}
% ──────────────────────────────────────────────

\subsection{Definitions and Examples}

\begin{definition}[Limit Cycle]
A \textbf{limit cycle} is an isolated closed orbit.  ``Isolated'' means
that neighboring trajectories are not closed---they either spiral toward or
away from the limit cycle.
\begin{itemize}[leftmargin=2em]
\item \textbf{Stable} limit cycle: all nearby trajectories approach it as
      $t\to\infty$.
\item \textbf{Unstable} limit cycle: all nearby trajectories move away.
\item \textbf{Half-stable}: trajectories approach from one side and recede
      from the other.
\end{itemize}
\end{definition}

\begin{remark}
Limit cycles are inherently \textbf{nonlinear} phenomena.  They cannot
occur in linear systems (where closed orbits, if they exist, form
continuous families---centers---not isolated cycles).
\end{remark}

\begin{example}[Limit Cycle in Polar Coordinates]
Consider $\dot{r} = r(1-r^2)$, $\dot{\theta} = 1$.
The radial equation has fixed points $r^*=0$ (unstable, since
$f'(0)=1>0$) and $r^*=1$ (stable, since $f'(1)=-2<0$).  Therefore
$r=1$ is a \textbf{stable limit cycle}: all trajectories with $r>0$
spiral toward the unit circle.
\end{example}

\subsection{Ruling Out Closed Orbits}

Before searching for limit cycles, it is often useful to rule them out in
certain regions.

\subsubsection{Gradient Systems}

If $\dot{x} = -\nabla V$ for some potential $V(x,y)$, then
\[
  \dv{V}{t} = \nabla V \cdot \dot{\vec{x}}
  = \nabla V \cdot (-\nabla V) = -|\nabla V|^2 \le 0.
\]
So $V$ decreases monotonically along trajectories, and closed orbits are
impossible.

\subsubsection{Liapunov Functions}

\begin{theorem}[Liapunov's Theorem for Ruling Out Closed Orbits]
If there exists a function $V(x,y)$ satisfying:
\begin{enumerate}[leftmargin=2em]
\item $V(\xstar,y^*)=0$ at all fixed points,
\item $V(x,y)>0$ for all $(x,y)$ not a fixed point (in some region $R$),
\item $\dot{V} = \dv{V}{t} < 0$ for all non-fixed points in $R$,
\end{enumerate}
then there are no closed orbits entirely contained in $R$.
\end{theorem}

\begin{proof}[Proof sketch]
On a closed orbit, $V$ would have to return to its starting value after one
period.  But $\dot{V}<0$ means $V$ is strictly decreasing, a contradiction.
\end{proof}

\subsubsection{Dulac's Criterion}

\begin{theorem}[Dulac's Criterion]\label{thm:dulac}
Let $R$ be a simply connected region of the plane.  If there exists a
continuously differentiable function $g(x,y)$ such that
$\nabla\cdot(g\vec{F})$ does not change sign in $R$ (and is not
identically zero on any open subset), then there are no closed orbits
entirely contained in $R$.

Here $\vec{F} = (f, g_{\text{sys}})$ is the vector field, and we require
\begin{equation}
  \pdv{(g\,f)}{x} + \pdv{(g\,g_{\text{sys}})}{y} \neq 0
  \quad\text{(same sign throughout $R$)}.
\end{equation}
\end{theorem}
\begin{proof}
By Green's theorem, for any closed curve $C$ bounding a region $A\subset R$:
\[
  \iint_A \nabla\cdot(g\vec{F})\,dA = \oint_C g\vec{F}\cdot\hat{n}\,dl.
\]
If $C$ were a closed orbit, then $\vec{F}$ is tangent to $C$ everywhere, so
$\vec{F}\cdot\hat{n}=0$ and the right side vanishes.  But the left side is
nonzero by hypothesis, a contradiction.
\end{proof}

\subsection{Poincar\'e--Bendixson Theorem}

This theorem is one of the most powerful tools for \emph{proving the
existence} of limit cycles.

\begin{theorem}[Poincar\'e--Bendixson Theorem]\label{thm:PB}
Suppose that:
\begin{enumerate}[leftmargin=2em]
\item $R$ is a closed, bounded region of the plane;
\item $\dot{\vec{x}} = \vec{F}(\vec{x})$ is a continuously differentiable
      vector field on an open set containing $R$;
\item $R$ does not contain any fixed points;
\item There exists a trajectory $C$ that is ``confined'' to $R$
      (i.e., $C$ starts in $R$ and can never leave).
\end{enumerate}
Then $C$ is either a closed orbit, or it spirals toward a closed orbit
contained in $R$.  In either case, $R$ contains a \textbf{closed orbit}.
\end{theorem}

\begin{remark}
This theorem is specific to two dimensions.  In three or more dimensions,
a bounded trajectory need not approach a closed orbit (it could be a
strange attractor, for example).
\end{remark}

\begin{physicalnote}[Trapping Regions]
To apply the Poincar\'e--Bendixson theorem, one constructs a
\textbf{trapping region}: a closed, bounded region $R$ with no fixed points
inside, such that the vector field points inward (or is tangent) on the
entire boundary $\partial R$.  Then any trajectory entering $R$ is confined
to it, and the theorem guarantees a closed orbit inside.

The boundary of $R$ is often built from segments of nullclines, coordinate
axes, or level curves of auxiliary functions.
\end{physicalnote}

\subsection{Li\'enard Systems}\label{sec:lienard}

\begin{definition}[Li\'enard System]
A \textbf{Li\'enard system} is a second-order ODE of the form
\begin{equation}
  \ddot{x} + f(x)\dot{x} + g(x) = 0,
\end{equation}
or equivalently, the system $\dot{x} = y - F(x)$, $\dot{y} = -g(x)$,
where $F(x) = \int_0^x f(s)\,ds$.
\end{definition}

\begin{theorem}[Li\'enard's Theorem]
The Li\'enard system has a unique, stable limit cycle if:
\begin{enumerate}[leftmargin=2em]
\item $f$ and $g$ are continuous, and $g$ is Lipschitz;
\item $g$ is an odd function: $g(-x) = -g(x)$, with $g(x) > 0$ for $x > 0$;
\item $F$ is an odd function with exactly one positive zero at $x = a$,
      $F(x) < 0$ for $0 < x < a$, $F(x) > 0$ and nondecreasing for
      $x > a$, and $F(x) \to \infty$ as $x \to \infty$.
\end{enumerate}
\end{theorem}

\subsection{The Van der Pol Oscillator}\label{sec:vdp}

The most important Li\'enard system is the \textbf{Van der Pol oscillator}:
\begin{equation}\label{eq:vdp}
  \ddot{x} + \mu(x^2-1)\dot{x} + x = 0, \qquad \mu > 0.
\end{equation}
Here $f(x) = \mu(x^2-1)$ and $g(x) = x$.  The damping is
\emph{negative} for $|x|<1$ (energy is pumped in) and \emph{positive}
for $|x|>1$ (energy is dissipated).  This self-regulating mechanism
produces a unique, stable limit cycle for any $\mu > 0$.

Rewriting as a system with $y = \dot{x}/\mu + x^3/3 - x$ (Li\'enard
plane):
\begin{equation}
  \dot{x} = \mu\bigl(y - F(x)\bigr), \qquad \dot{y} = -\frac{x}{\mu},
  \qquad F(x) = \frac{x^3}{3} - x.
\end{equation}

\subsection{Relaxation Oscillations}\label{sec:relaxation}

For $\mu \gg 1$ (strong nonlinearity), the Van der Pol oscillator exhibits
\textbf{relaxation oscillations}: the motion consists of slow crawls along
the cubic nullcline $y = F(x) = x^3/3 - x$ punctuated by fast jumps
between the two branches.

\textbf{Mechanism:}
\begin{enumerate}[leftmargin=2em]
\item \textbf{Fast dynamics} ($\dot{x}$ is large): the trajectory jumps
      nearly horizontally between the left and right branches of the
      $S$-shaped nullcline.
\item \textbf{Slow dynamics} ($\dot{y}$ is small): the trajectory crawls
      along the nullcline until it reaches a knee (fold), then jumps again.
\end{enumerate}

The period of relaxation oscillations scales as $T \approx (3 - 2\ln 2)\mu
\approx 1.614\,\mu$ for large $\mu$---it grows linearly with the
nonlinearity parameter, in contrast to the nearly constant period of
weakly nonlinear oscillations.


% ──────────────────────────────────────────────
\section{Bifurcations Revisited: Two-Dimensional Systems}%
\label{sec:bif-2d}
% ──────────────────────────────────────────────

\subsection{Saddle-Node, Transcritical, and Pitchfork in 2D}

In two dimensions, all three bifurcation types from Section~\ref{sec:bifurcations}
can occur.  The mechanism is the same: one eigenvalue of the Jacobian passes
through zero.  The other eigenvalue remains nonzero, so the dynamics in the
``transverse'' direction is simply attracted to (or repelled from) a
manifold on which the 1D bifurcation takes place.

\subsection{Hopf Bifurcation}\label{sec:hopf}

The \textbf{Hopf bifurcation} is the birth (or death) of a limit cycle from
a fixed point.  It occurs when a pair of complex conjugate eigenvalues
crosses the imaginary axis.

\begin{theorem}[Hopf Bifurcation Theorem---Informal]
Consider $\dot{\vec{x}} = \vec{F}(\vec{x};\mu)$ with a fixed point at
the origin.  Suppose the Jacobian at the origin has eigenvalues
$\lambda(\mu) = \alpha(\mu) \pm i\omega(\mu)$.  If:
\begin{enumerate}[leftmargin=2em]
\item At $\mu=0$: $\alpha(0) = 0$ and $\omega(0) \neq 0$ (pure imaginary
      eigenvalues);
\item $\alpha'(0) \neq 0$ (eigenvalues cross the imaginary axis with
      nonzero speed);
\end{enumerate}
then a Hopf bifurcation occurs at $\mu=0$.
\end{theorem}

\subsubsection{Supercritical Hopf Bifurcation}

A \textbf{stable limit cycle} is born from the fixed point as it loses
stability.

\begin{keyresult}[Normal Form: Supercritical Hopf]
In polar coordinates:
\begin{equation}\label{eq:hopf-super}
  \dot{r} = \mu r - r^3, \qquad \dot{\theta} = \omega + b r^2.
\end{equation}
\begin{itemize}[leftmargin=2em]
\item For $\mu < 0$: origin is a stable spiral.
\item For $\mu > 0$: origin becomes unstable; a stable limit cycle appears
      at $r = \sqrt{\mu}$.
\end{itemize}
The amplitude of the limit cycle grows as $\sqrt{\mu}$, and its frequency
is $\omega + b\mu$ (the parameter $b$ shifts the frequency, while
$\omega$ sets the base frequency).
\end{keyresult}

\subsubsection{Subcritical Hopf Bifurcation}

An \textbf{unstable} limit cycle shrinks down to the origin and is absorbed
as the origin gains stability.

Normal form in polar coordinates:
\begin{equation}
  \dot{r} = \mu r + r^3, \qquad \dot{\theta} = \omega + b r^2.
\end{equation}
\begin{itemize}[leftmargin=2em]
\item For $\mu < 0$: an unstable limit cycle at $r=\sqrt{-\mu}$ coexists
      with a stable origin.
\item At $\mu = 0$: the limit cycle collapses into the origin.
\item For $\mu > 0$: the origin is unstable and no small-amplitude limit
      cycle exists.  Trajectories may jump to a distant attractor, leading
      to a \textbf{hard} or \textbf{catastrophic} transition.
\end{itemize}

\subsection{Global Bifurcations of Cycles}

Limit cycles can also be born or destroyed through \emph{global}
bifurcations that involve large-scale rearrangements of the phase portrait.

\begin{enumerate}[leftmargin=2em]
\item \textbf{Saddle-node bifurcation of cycles (fold of cycles):}
      A stable and an unstable limit cycle collide and annihilate.
\item \textbf{Homoclinic bifurcation:} A limit cycle grows until it
      touches a saddle point, forming a \emph{homoclinic orbit} (a
      trajectory that connects the saddle to itself).  As the parameter
      passes through the bifurcation, the cycle's period diverges
      ($T\to\infty$), and the cycle disappears.
\item \textbf{Infinite-period (SNIPER) bifurcation:} A limit cycle passes
      through a saddle-node on a circle.  The period diverges but the
      amplitude remains finite.
\end{enumerate}

\subsection{Hysteresis in the Josephson Junction (2D)}\label{sec:jj-2d}

With capacitance included, the Josephson junction becomes a second-order
system:
\begin{equation}
  i = \sin\phi + \alpha\dot{\phi} + \ddot{\phi},
\end{equation}
which is equivalent to the \textbf{damped driven pendulum}:
\begin{equation}
  \dot{\phi} = y, \qquad \dot{y} = i - \sin\phi - \alpha y.
\end{equation}
The Jacobian is
$J = \bigl(\begin{smallmatrix} 0 & 1 \\ -\cos\phi & -\alpha
\end{smallmatrix}\bigr)$.

The parameter $\alpha$ controls the damping:
\begin{itemize}[leftmargin=2em]
\item \textbf{Overdamped} ($\alpha \gg 1$): reduces to the 1D model of
      Section~\ref{sec:josephson-1d}.
\item \textbf{Underdamped} ($\alpha < 1$): hysteresis occurs.  As $i$
      increases past $1$, the zero-voltage state is lost; but when $i$ is
      decreased, the system remains on the rotating (voltage) branch until
      $i$ drops below a retrapping current $i_r < 1$.
\end{itemize}

\subsection{Coupled Oscillators and Quasiperiodicity}\label{sec:coupled-osc}

Consider two oscillators with phase variables $\theta_1,\theta_2$ on the
circle, coupled weakly:
\begin{equation}\label{eq:coupled-osc}
  \dot{\theta}_1 = \omega_1 + K_1 \sin(\theta_2 - \theta_1), \qquad
  \dot{\theta}_2 = \omega_2 + K_2 \sin(\theta_1 - \theta_2).
\end{equation}
Defining the phase difference $\phi = \theta_1 - \theta_2$:
\begin{equation}
  \dot{\phi} = (\omega_1 - \omega_2) - (K_1 + K_2)\sin\phi.
\end{equation}
This is again a flow on the circle!

\begin{itemize}[leftmargin=2em]
\item \textbf{Phase locking:}  If $|\omega_1 - \omega_2| < K_1 + K_2$,
      stable fixed points exist: the oscillators lock to a constant phase
      difference.
\item \textbf{Phase drift:}  If $|\omega_1 - \omega_2| > K_1 + K_2$,
      no fixed points; $\phi$ drifts and the oscillators are
      \textbf{not synchronized}.
\end{itemize}

\subsubsection{Dynamics on the Torus}

The combined phase space $(\theta_1,\theta_2)$ lives on a \textbf{torus}
$\mathbb{T}^2 = S^1 \times S^1$.  When uncoupled ($K_1=K_2=0$), each
oscillator moves independently, and the trajectory on the torus is either:
\begin{itemize}[leftmargin=2em]
\item A \textbf{closed orbit} if $\omega_1/\omega_2$ is rational, or
\item A \textbf{quasiperiodic orbit} (densely filling the torus) if
      $\omega_1/\omega_2$ is irrational.
\end{itemize}

\subsection{Poincar\'e Maps}\label{sec:poincare-maps}

\begin{definition}[Poincar\'e Map]
Choose a \textbf{section} $\Sigma$ (a codimension-1 surface transverse to
the flow) in the phase space.  The \textbf{Poincar\'e map} (or first-return
map) $P:\Sigma\to\Sigma$ sends each point $\vec{x}\in\Sigma$ to the next
point where the trajectory through $\vec{x}$ intersects $\Sigma$.
\end{definition}

This reduces a continuous-time system to a discrete-time map, lowering the
dimension by one.  A \textbf{closed orbit} of the flow corresponds to a
\textbf{fixed point} of the Poincar\'e map.


% ══════════════════════════════════════════════
\part{Chaos}
% ══════════════════════════════════════════════

% ──────────────────────────────────────────────
\section{The Lorenz Equations}\label{sec:lorenz}
% ──────────────────────────────────────────────

\subsection{Origin and Physical Motivation}

Edward Lorenz (1963) derived a simplified model of atmospheric convection
(Rayleigh--B\'enard convection).  The system consists of three coupled ODEs:
\begin{equation}\label{eq:lorenz}
\boxed{
  \dot{x} = \sigma(y-x), \qquad
  \dot{y} = rx - y - xz, \qquad
  \dot{z} = xy - bz,
}
\end{equation}
where $\sigma > 0$ is the Prandtl number, $r > 0$ is (proportional to) the
Rayleigh number, and $b > 0$ is a geometric factor.  All parameters are
positive.  The classic parameter values are $\sigma = 10$, $b = 8/3$.

\subsection{Simple Properties}

\begin{enumerate}[leftmargin=2em]
\item \textbf{Symmetry:}  The system is invariant under
      $(x,y,z)\to(-x,-y,z)$.

\item \textbf{Dissipation:}  The divergence of the flow is
      \[
        \nabla\cdot\vec{F} = -\sigma - 1 - b < 0,
      \]
      so volumes in phase space contract exponentially.  This means the
      system is \textbf{dissipative} and all trajectories eventually
      become confined to an attractor of zero volume.

\item \textbf{Fixed points:}
      \begin{itemize}
      \item The origin $O = (0,0,0)$ exists for all $r$.
      \item For $r > 1$, two additional fixed points appear via a
            \textbf{supercritical pitchfork} bifurcation:
            \[
              C^\pm = \bigl(\pm\sqrt{b(r-1)},\;\pm\sqrt{b(r-1)},\;r-1\bigr).
            \]
      \end{itemize}

\item \textbf{Stability of the origin:}  The eigenvalues at $O$ are found
      from the Jacobian.  The origin is stable for $r<1$ and unstable (a
      saddle with one unstable direction) for $r>1$.

\item \textbf{Stability of $C^\pm$:}  These are stable spirals for
      $1 < r < r_H$, where
      \begin{equation}
        r_H = \sigma\,\frac{\sigma + b + 3}{\sigma - b - 1}
      \end{equation}
      (assuming $\sigma > b+1$).  At $r=r_H$, a \textbf{subcritical Hopf
      bifurcation} occurs: $C^\pm$ lose stability and an unstable limit
      cycle is absorbed.

\item \textbf{Trapping region:}  Using the Liapunov-like function
      $V = rx^2 + \sigma y^2 + \sigma(z-2r)^2$, one can show $\dot{V}<0$
      outside a large ellipsoid.  Hence all trajectories eventually enter
      this region and remain bounded.
\end{enumerate}

\subsection{Chaos on the Strange Attractor}

For $r > r_H$ (and the classic parameters), trajectories exhibit
\textbf{chaotic} behavior:
\begin{itemize}[leftmargin=2em]
\item \textbf{Aperiodic:} Trajectories never repeat, even though they are
      bounded.
\item \textbf{Sensitive dependence on initial conditions:}  Nearby
      trajectories diverge exponentially.  This is the hallmark of chaos.
\item \textbf{Strange attractor:}  The attractor has a fractal structure
      with zero volume but infinite detail.
\end{itemize}

\begin{definition}[Sensitive Dependence and Horizon Time]
Two nearby initial conditions separated by $\delta_0$ diverge roughly as
$\delta(t) \sim \delta_0\, e^{\lambda t}$, where $\lambda > 0$ is the
\textbf{Lyapunov exponent}.  Predictions become unreliable when
$\delta(t) \sim \epsilon$ (some tolerance).  The \textbf{horizon time} is:
\begin{equation}\label{eq:horizon}
  t_h \sim \frac{1}{\lambda}\ln\frac{\epsilon}{\delta_0}.
\end{equation}
Beyond $t_h$, prediction is essentially impossible despite deterministic
dynamics.  Note the logarithmic dependence: even dramatic improvements in
measurement precision $\delta_0$ yield only modest increases in $t_h$.
\end{definition}

\subsection{The Lorenz Map}

Lorenz discovered that the chaotic dynamics can be partially understood by
a one-dimensional map.  Record the successive local maxima $z_n$ of $z(t)$.
Plotting $z_{n+1}$ vs.\ $z_n$ yields a remarkably thin, single-humped curve
(approximately a tent map).

Key observations:
\begin{itemize}[leftmargin=2em]
\item The slope $|f'| > 1$ everywhere on the Lorenz map, implying that
      nearby orbits are stretched apart at each return.
\item There are no stable fixed points or periodic orbits of the map,
      explaining the aperiodic behavior.
\item The thinness of the map (nearly 1D) arises from the extreme
      contraction in the dissipative system.
\end{itemize}

\subsection{Exploring Parameter Space of the Lorenz System}

As $r$ increases (with $\sigma=10$, $b=8/3$), the Lorenz system passes
through several regimes:
\begin{enumerate}[leftmargin=2em]
\item $0 < r < 1$: The origin is the global attractor.  All motion decays.
\item $1 < r < r_H \approx 24.74$: The fixed points $C^\pm$ are stable
      spirals.  Transient chaos may occur for $r > 13.9$ (where a
      ``strange saddle'' exists), but all trajectories eventually settle
      to $C^+$ or $C^-$.
\item $r_H < r$: Chaotic attractor (the ``butterfly'').  The system
      exhibits sustained aperiodic motion, sensitive dependence, and a
      fractal attractor.
\item Periodic windows exist within the chaotic regime, where stable
      periodic orbits briefly emerge.
\end{enumerate}

\subsection{The AC-Driven Pendulum and Routes to Chaos}\label{sec:ac-pendulum}

The damped pendulum with a periodic (AC) driving torque provides another
canonical route to chaos:
\begin{equation}
  \Gamma\cos(\omega t) = \ddot{\theta} + b\dot{\theta} + \sin\theta.
\end{equation}
Writing as a system with $\dot{\theta} = v$, $\dot{v} = -bv - \sin\theta +
\Gamma\cos(\omega y)$, $\dot{y} = 1$, this is a three-dimensional
autonomous system (the explicit time dependence elevates it to 3D).

The Poincar\'e section (strobing at the drive frequency) reduces it to a
2D map.  As $\Gamma$ increases, the system passes through:
\begin{itemize}[leftmargin=2em]
\item Periodic (phase-locked) motion,
\item Period-doubling cascades,
\item Chaotic tumbling,
\end{itemize}
providing a concrete physical system exhibiting the period-doubling route
to chaos.


% ──────────────────────────────────────────────
\section{What Is Chaos?}\label{sec:chaos-defn}
% ──────────────────────────────────────────────

Before proceeding to maps, let us formalize the notion of chaos.

\begin{definition}[Chaos---Working Definition]
A dynamical system is \textbf{chaotic} if it satisfies all three conditions:
\begin{enumerate}[leftmargin=2em]
\item \textbf{Aperiodic long-term behavior:} Trajectories do not settle
      to fixed points, periodic orbits, or quasiperiodic orbits.
\item \textbf{Deterministic:} There is no randomness in the equations---the
      system is governed by definite rules with no stochastic terms.
\item \textbf{Sensitive dependence on initial conditions:} Nearby
      trajectories diverge exponentially, with at least one positive
      Lyapunov exponent.
\end{enumerate}
\end{definition}

\begin{remark}
Chaos is \emph{not} randomness.  The underlying equations are perfectly
deterministic; the apparent randomness arises from the amplification of
microscopic uncertainties in initial conditions.  This is sometimes called
\textbf{deterministic chaos} or \textbf{deterministic aperiodicity}.
\end{remark}

\begin{remark}
Chaos requires at least three dimensions for continuous-time (flow) systems.
In two dimensions, the Poincar\'e--Bendixson theorem constrains bounded
trajectories to approach fixed points or closed orbits.  However, for
discrete-time systems (maps), chaos is possible even in one dimension
(e.g., the logistic map).
\end{remark}

% ──────────────────────────────────────────────
\section{One-Dimensional Maps}\label{sec:1d-maps}
% ──────────────────────────────────────────────

\subsection{Introduction}

Motivated by Poincar\'e maps and the Lorenz map, we now study
discrete-time dynamical systems (\textbf{iterated maps}) of the form
\begin{equation}\label{eq:map}
  x_{n+1} = f(x_n),
\end{equation}
where $f:\R\to\R$ is a given function.

\subsection{Fixed Points and Cobweb Diagrams}

\begin{definition}
A point $\xstar$ is a \textbf{fixed point} of the map~\eqref{eq:map} if
$f(\xstar) = \xstar$.
\end{definition}

Graphically, fixed points are found by intersecting the graph of $f(x)$
with the diagonal line $y=x$.

\begin{keyresult}[Stability of a Fixed Point (Maps)]
Let $\xstar$ be a fixed point of $x_{n+1}=f(x_n)$.  The \textbf{multiplier}
is $\mu = f'(\xstar)$.
\begin{itemize}[leftmargin=2em]
\item $|\mu| < 1$: $\xstar$ is \textbf{stable} (attracting).
\item $|\mu| > 1$: $\xstar$ is \textbf{unstable} (repelling).
\item $|\mu| = 1$: inconclusive.
\end{itemize}
\end{keyresult}

\begin{proof}
Let $\delta_n = x_n - \xstar$.  Then
$\delta_{n+1} = f(\xstar + \delta_n) - \xstar \approx f'(\xstar)\,\delta_n$,
so $\delta_n \approx [f'(\xstar)]^n \delta_0$.
\end{proof}

\textbf{Cobweb construction:}  A graphical method for iterating maps.
Starting from $(x_0, 0)$:
\begin{enumerate}[leftmargin=2em]
\item Go vertically to $y = f(x_0)$.
\item Go horizontally to the diagonal $y = x$ (so the horizontal
      coordinate is now $x_1 = f(x_0)$).
\item Repeat.
\end{enumerate}

\subsection{Periodic Orbits}

\begin{definition}
A point $p$ has \textbf{period $n$} if $f^n(p) = p$ and $f^k(p)\neq p$
for $0 < k < n$, where $f^n$ denotes the $n$-fold composition
$f\circ f\circ\cdots\circ f$.
\end{definition}

A period-$n$ orbit is a fixed point of $f^n$.  Its stability is determined
by the multiplier of $f^n$:
\begin{equation}
  (f^n)'(p) = f'(p)\cdot f'(f(p))\cdot f'(f^2(p))\cdots f'(f^{n-1}(p))
  = \prod_{k=0}^{n-1} f'(f^k(p)).
\end{equation}
The orbit is stable if $|(f^n)'(p)| < 1$ and unstable if $|(f^n)'(p)| > 1$.

\subsection{The Logistic Map}\label{sec:logistic}

The most studied one-dimensional map is the \textbf{logistic map}:
\begin{equation}\label{eq:logistic}
  x_{n+1} = rx_n(1-x_n), \qquad 0 \le r \le 4.
\end{equation}

\subsubsection{Fixed Points}

Setting $x = rx(1-x)$: the fixed points are $\xstar = 0$ and
$\xstar = (r-1)/r$ (for $r>1$).

Stability of $\xstar = (r-1)/r$:
$f'(x) = r - 2rx$, so $f'(\xstar) = r - 2(r-1) = 2 - r$.
This fixed point is stable when $|2-r|<1$, i.e., $1 < r < 3$.

\subsubsection{Period Doubling}

At $r = 3$, the fixed point $\xstar = (r-1)/r$ loses stability through a
\textbf{flip bifurcation} ($f'(\xstar) = -1$).  A stable period-2 cycle is
born.  As $r$ increases further:
\begin{itemize}[leftmargin=2em]
\item $r \approx 3$: period-2 cycle born.
\item $r \approx 3.449$: period-4 cycle born.
\item $r \approx 3.544$: period-8 cycle born.
\item The cascade continues: period $2^n$ orbits for $n = 1, 2, 3, \ldots$
\item The intervals between successive doublings shrink geometrically.
\item At $r_\infty \approx 3.5699$: all periods $2^n$ coexist---the onset
      of chaos.
\end{itemize}

\subsubsection{Period-Doubling Cascade and Feigenbaum's Universality}

Let $r_n$ be the parameter value at which the period-$2^n$ cycle is born.
Feigenbaum (1978) discovered that the ratios
\begin{equation}
  \frac{r_n - r_{n-1}}{r_{n+1} - r_n} \to \delta \approx 4.669\,201\ldots
\end{equation}
converge to a universal constant $\delta$ (\textbf{Feigenbaum's constant}),
independent of the specific map $f$, as long as $f$ has a single quadratic
maximum.

\subsubsection{Periodic Windows}

For $r > r_\infty$, chaos is interspersed with \textbf{periodic windows}---
intervals of $r$ where a stable periodic orbit exists.  The most prominent
is the period-3 window near $r \approx 3.83$.

\begin{theorem}[Li--Yorke Theorem (1975)]
If a continuous map on an interval has a period-3 orbit, then it has
periodic orbits of \emph{every} period.
\end{theorem}

\subsection{Lyapunov Exponent}\label{sec:lyapunov-exp}

\begin{definition}[Lyapunov Exponent for Maps]
For a map $x_{n+1} = f(x_n)$ with initial condition $x_0$, the
\textbf{Lyapunov exponent} is
\begin{equation}\label{eq:lyap-map}
  \lambda = \lim_{n\to\infty}\frac{1}{n}\sum_{k=0}^{n-1}\ln|f'(x_k)|
  = \lim_{n\to\infty}\frac{1}{n}\ln\left|\frac{d f^n}{dx}\bigg|_{x_0}\right|.
\end{equation}
\end{definition}

\textbf{Interpretation:}
\begin{itemize}[leftmargin=2em]
\item $\lambda > 0$: nearby orbits diverge exponentially---the system is
      \textbf{chaotic}.
\item $\lambda < 0$: nearby orbits converge---a stable periodic orbit
      attracts.
\item $\lambda = 0$: marginal case (e.g., at a bifurcation point or
      neutral stability).
\end{itemize}

For two nearby initial conditions $x_0$ and $\tilde{x}_0 = x_0 + \delta_0$:
\[
  |\delta_n| \approx |\delta_0|\,e^{\lambda n}
  \quad\Longrightarrow\quad
  \lambda \approx \frac{1}{n}\ln\frac{|\delta_n|}{|\delta_0|}.
\]

\subsection{Renormalization}\label{sec:renormalization}

The universality of the period-doubling cascade can be understood through
\textbf{renormalization group} ideas.

The key insight: near the period-doubling accumulation point $r_\infty$,
the map $f^{2^n}$ (restricted to a small interval near the critical point)
looks approximately like a rescaled copy of $f$ itself.  This self-similarity
leads to a functional equation (the \textbf{Feigenbaum equation}):
\begin{equation}
  g(x) = -\alpha\, g\bigl(g(-x/\alpha)\bigr),
\end{equation}
where $g$ is the universal limiting function and
$\alpha \approx 2.5029$ is the rescaling factor.  The constant $\delta$
(Feigenbaum's constant) arises as the relevant eigenvalue of the
renormalization operator linearized about $g$.


% ──────────────────────────────────────────────
\section{Fractals and Strange Attractors}\label{sec:fractals}
% ──────────────────────────────────────────────

\subsection{Fractal Dimension}

Chaotic attractors typically have a \textbf{fractal} (non-integer) dimension.
This reflects their infinitely detailed, self-similar structure.

\begin{definition}[Box-Counting Dimension]
Cover the set $S$ with $N(\epsilon)$ boxes of side length $\epsilon$.  The
\textbf{box-counting dimension} (or Minkowski dimension) is
\begin{equation}
  d = \lim_{\epsilon\to 0}\frac{\ln N(\epsilon)}{\ln(1/\epsilon)}.
\end{equation}
\end{definition}

\begin{example}[Self-Similar Fractals]
\begin{itemize}[leftmargin=2em]
\item \textbf{Cantor set:} Formed by repeatedly removing the middle third
      of intervals.  Dimension $d = \ln 2/\ln 3 \approx 0.631$.
\item \textbf{Sierpinski triangle:} Dimension $d = \ln 3/\ln 2 \approx 1.585$.
\item \textbf{Koch snowflake:} Dimension $d = \ln 4/\ln 3 \approx 1.262$.
\end{itemize}
For a self-similar fractal made of $N$ copies of itself scaled by factor
$r$, the dimension is $d = \ln N / \ln(1/r)$.
\end{example}

\subsection{Dimension of Self-Similar Fractals}

\begin{definition}[Self-Similarity Dimension]
If a set $S$ is the union of $N$ non-overlapping copies of itself, each
scaled by a factor $r$ (with $0 < r < 1$), then the
\textbf{self-similarity dimension} is
\begin{equation}
  d = \frac{\ln N}{\ln(1/r)}.
\end{equation}
\end{definition}

This coincides with the box-counting dimension for exact self-similar sets.
For strange attractors (which are typically self-similar only in a
statistical sense), one uses the box-counting or correlation dimension.

\subsection{Strange Attractors}

\begin{definition}[Strange Attractor]
An attractor is \textbf{strange} if it has a fractal structure (non-integer
dimension).  A \textbf{chaotic attractor} is a strange attractor on which
trajectories exhibit sensitive dependence on initial conditions.
\end{definition}

Key examples:
\begin{itemize}[leftmargin=2em]
\item The \textbf{Lorenz attractor} has dimension $\approx 2.06$---slightly
      more than a surface, reflecting its layered structure.
\item The \textbf{H\'enon map} $x_{n+1} = 1 - ax_n^2 + y_n$,
      $y_{n+1} = bx_n$ (with $a=1.4$, $b=0.3$) has a strange attractor
      with dimension $\approx 1.26$.
\item The \textbf{R\"ossler system}
      $\dot{x} = -y-z$, $\dot{y} = x+ay$, $\dot{z} = b+z(x-c)$ produces
      a simpler ``band'' attractor that is easier to visualize than Lorenz.
\end{itemize}

\begin{physicalnote}[Stretching and Folding]
Strange attractors arise from the interplay of two mechanisms:
\begin{enumerate}[leftmargin=2em]
\item \textbf{Stretching:} Nearby trajectories diverge (positive Lyapunov
      exponent), creating sensitive dependence.
\item \textbf{Folding:} Since the attractor is bounded, the stretched-out
      trajectories must fold back on themselves, creating the layered
      fractal structure.
\end{enumerate}
This stretch-and-fold mechanism is analogous to kneading dough---each
iteration creates finer and finer layers.
\end{physicalnote}


% ══════════════════════════════════════════════
\part{Collective Behavior}
% ══════════════════════════════════════════════

% ──────────────────────────────────────────────
\section{The Kuramoto Model}\label{sec:kuramoto}
% ──────────────────────────────────────────────

\subsection{Motivation: Synchronization in Nature}

Synchronization phenomena are ubiquitous: fireflies flashing in unison,
cardiac pacemaker cells, audience clapping in rhythm, circadian rhythms, and
power grids.  The challenge is to understand how a large population of
coupled oscillators can spontaneously synchronize.

\subsection{The Model}

Yoshiki Kuramoto (1975) proposed a beautifully simple model.  Consider $N$
oscillators with phases $\theta_1, \ldots, \theta_N$.  Each oscillator has
its own natural frequency $\omega_i$, drawn from a distribution $g(\omega)$
(typically symmetric and unimodal, e.g., Gaussian or Lorentzian).  The
coupling is all-to-all and sinusoidal:

\begin{equation}\label{eq:kuramoto}
\boxed{
  \dot{\theta}_i = \omega_i + \frac{K}{N}\sum_{j=1}^{N}
  \sin(\theta_j - \theta_i), \qquad i = 1,\ldots,N.
}
\end{equation}

Here $K \ge 0$ is the \textbf{coupling strength}.

\subsection{The Order Parameter}

To measure the degree of synchronization, Kuramoto introduced the
\textbf{complex order parameter}:
\begin{equation}\label{eq:order-param}
  re^{i\psi} = \frac{1}{N}\sum_{j=1}^{N} e^{i\theta_j},
\end{equation}
where $r\in[0,1]$ is the \textbf{coherence} (amplitude of the mean field)
and $\psi$ is the \textbf{average phase}.

\begin{itemize}[leftmargin=2em]
\item $r = 1$: all oscillators are perfectly synchronized
      ($\theta_j = \psi$ for all $j$).
\item $r = 0$: the phases are uniformly distributed around the circle
      (complete incoherence).
\item $0 < r < 1$: partial synchronization.
\end{itemize}

\textbf{Mean-field form.}  Multiplying both sides of~\eqref{eq:order-param}
by $e^{-i\theta_i}$ and taking the imaginary part:
\[
  r\sin(\psi - \theta_i) = \frac{1}{N}\sum_{j=1}^{N}
  \sin(\theta_j - \theta_i).
\]
Thus the Kuramoto model becomes:
\begin{equation}\label{eq:kuramoto-mf}
  \dot{\theta}_i = \omega_i + Kr\sin(\psi - \theta_i).
\end{equation}
Each oscillator interacts only with the ``mean field'' $(r,\psi)$, not with
every other oscillator individually.

\subsection{Self-Consistency Analysis (Infinite $N$)}

In the $N\to\infty$ limit, we look for a steady-state solution in a
rotating frame where $\psi$ is constant (choose a frame where $\psi=0$).
Oscillators split into two groups:

\begin{enumerate}[leftmargin=2em]
\item \textbf{Locked oscillators} ($|\omega_i| \le Kr$): these are
      entrained to the mean field and satisfy $\dot{\theta}_i = 0$, i.e.,
      $\omega_i = -Kr\sin\theta_i$, giving
      $\theta_i = \arcsin(\omega_i / Kr)$.

\item \textbf{Drifting oscillators} ($|\omega_i| > Kr$): these are not
      captured by the mean field and continue to rotate, but nonuniformly.
      Their steady-state density is inversely proportional to their angular
      velocity: $\rho(\omega,\theta) \propto 1/|\dot{\theta}|$.
\end{enumerate}

The self-consistency condition requires the locked and drifting oscillators
to produce exactly the assumed order parameter $r$.  This yields:
\begin{equation}\label{eq:self-consistency}
  r = \int_{-\pi/2}^{\pi/2} Kr\cos^2\theta\; g(Kr\sin\theta)\,d\theta,
\end{equation}
or equivalently, after canceling $r$ (for $r\neq 0$):
\begin{equation}\label{eq:self-consistency-2}
  1 = K\int_{-\pi/2}^{\pi/2} \cos^2\theta\; g(Kr\sin\theta)\,d\theta.
\end{equation}

\subsection{Critical Coupling and the Phase Transition}

\begin{theorem}[Onset of Synchronization]
There exists a critical coupling strength
\begin{equation}\label{eq:Kc}
  K_c = \frac{2}{\pi\, g(0)},
\end{equation}
such that:
\begin{itemize}[leftmargin=2em]
\item For $K < K_c$: the only stable solution is $r = 0$
      (complete incoherence).
\item For $K > K_c$: a partially synchronized state with $r > 0$
      emerges.
\end{itemize}
Near the transition, $r \sim \sqrt{K - K_c}$, analogous to a
\textbf{supercritical pitchfork bifurcation} (or a second-order phase
transition in statistical mechanics).
\end{theorem}

\begin{proof}[Derivation of $K_c$]
From~\eqref{eq:self-consistency-2}, in the limit $r\to 0^+$:
\[
  1 = K\int_{-\pi/2}^{\pi/2}\cos^2\theta\;g(0)\,d\theta
  = Kg(0)\cdot\frac{\pi}{2},
\]
since $g(Kr\sin\theta)\to g(0)$ uniformly and
$\int_{-\pi/2}^{\pi/2}\cos^2\theta\,d\theta = \pi/2$.
Solving: $K_c = 2/(\pi g(0))$.
\end{proof}

\begin{physicalnote}[Analogy with Phase Transitions]
The Kuramoto model undergoes a transition remarkably similar to a
\textbf{ferromagnetic phase transition}: below $K_c$ (``high temperature''),
there is no order; above $K_c$ (``low temperature''), spontaneous order
emerges.  The order parameter $r$ plays the role of the magnetization, and
$K$ plays the role of inverse temperature.  The square-root scaling
$r\sim\sqrt{K-K_c}$ gives a critical exponent of $\beta = 1/2$, the same
as mean-field theory in statistical mechanics.
\end{physicalnote}


% ══════════════════════════════════════════════
% APPENDICES
% ══════════════════════════════════════════════
\appendix

\section{Summary of Key Theorems}\label{app:theorems}

\begin{enumerate}[leftmargin=2em]
\item \textbf{Existence and Uniqueness} (\ref{thm:exist-unique}):
      Guarantees local existence and uniqueness for $C^1$ vector fields.
\item \textbf{Hartman--Grobman} (\ref{thm:hartman-grobman}):
      Hyperbolic fixed points of nonlinear systems are topologically
      equivalent to the corresponding linear system.
\item \textbf{Dulac's Criterion} (\ref{thm:dulac}):
      Rules out closed orbits via divergence conditions.
\item \textbf{Poincar\'e--Bendixson} (\ref{thm:PB}):
      A confined trajectory in 2D must approach a closed orbit.
\item \textbf{Li--Yorke}:
      Period 3 implies chaos (existence of all periods).
\end{enumerate}

\section{Useful Mathematical Facts}\label{app:math-facts}

\subsection{Taylor Expansion for Stability Analysis}

For a function $f\in C^n$ near a point $\xstar$:
\[
  f(x) = f(\xstar) + f'(\xstar)(x-\xstar) + \frac{f''(\xstar)}{2}(x-\xstar)^2
  + \cdots + \frac{f^{(n)}(\xstar)}{n!}(x-\xstar)^n + O\bigl((x-\xstar)^{n+1}\bigr).
\]

\subsection{The Jacobian Matrix}

For a vector field $\dot{\vec{x}} = \vec{F}(\vec{x})$ in $\R^n$:
\[
  J_{ij} = \pdv{F_i}{x_j}\bigg|_{\vec{x}^*}.
\]
The eigenvalues of $J$ determine the local stability of the fixed point
$\vec{x}^*$.

\subsection{Green's Theorem (for Dulac's Criterion)}

For a continuously differentiable vector field
$\vec{A}(x,y) = (P(x,y),\,Q(x,y))$ and a simple closed curve $C$ bounding
a region $D$:
\[
  \oint_C (P\,dx + Q\,dy) = \iint_D \left(\pdv{Q}{x} - \pdv{P}{y}\right)dA.
\]

In divergence form (used for Dulac):
\[
  \iint_D \nabla\cdot\vec{A}\;dA = \oint_C \vec{A}\cdot\hat{n}\;dl.
\]

\subsection{Eigenvalues of a $2\times 2$ Matrix}

For $A = \bigl(\begin{smallmatrix} a & b \\ c & d \end{smallmatrix}\bigr)$:
\[
  \lambda = \frac{(a+d) \pm \sqrt{(a+d)^2 - 4(ad-bc)}}{2}
  = \frac{\tau \pm \sqrt{\tau^2 - 4\Delta}}{2}.
\]

\subsection{Polar Coordinate Conversions}

$x = r\cos\theta$, $y = r\sin\theta$; \quad
$r = \sqrt{x^2+y^2}$, $\theta = \arctan(y/x)$.

Useful derivatives:
\begin{align}
  r\dot{r} &= x\dot{x} + y\dot{y}, \\
  r^2\dot{\theta} &= x\dot{y} - y\dot{x}.
\end{align}

\end{document}
